% ============================================================
% Auto-generated from docs/golden-sunflowers/ch-27-tri27-dsl.md
% Source of truth: Railway phd-postgres-ssot ssot.chapters (gHashTag/trios#380)
% ============================================================

\chapter{TRI27 DSL}
\begin{figure}[H]
\centering
\makebox[\linewidth][c]{\includegraphics[width=1.18\linewidth,keepaspectratio]{ch27-tri27-dsl.png}}
\caption*{Figure — Ch.27: TRI27 DSL.}
\end{figure}

\begin{quote}
\itshape
``A language is a notation, and a notation is a thinking tool.''\\
\hfill --- Kenneth E.~Iverson, Turing Award Lecture, 1979.
\end{quote}

\section*{The smallest interesting alphabet}

Most languages have hundreds of letters or thousands of ideograms.
The smallest alphabet that still admits non-trivial logic is
\(\{-1, 0, +1\}\) --- three symbols, one of them neutral. TRI27 is the
domain-specific language built on this alphabet, and it is the kernel
language of the Trinity S\textsuperscript{3}AI system. Its three values
are not chosen for elegance; they are chosen because the integer
\emph{three} appears, with a Coq-checked proof, in the anchor identity
\(\varphi^{2} + \varphi^{-2} = 3\). Every program in TRI27 is
therefore a sequence of decisions over the same three-element
alphabet that the rest of this dissertation has been arguing for.

Two properties of the language are worth flagging at the door, before
the formal definition starts. First, evaluation is \emph{deterministic}:
the same expression in the same environment yields the same trit,
always, and we have a Coq theorem (\texttt{eval\_det}) that says so.
Second, the trit type is \emph{exhaustive}: every value of type
\texttt{trit} is one of exactly three possibilities, and we have a
second Coq theorem (\texttt{trit\_exhaustive}) that closes off any
attempt to smuggle in a fourth.

The rest of the chapter spells out the syntax, the denotational
semantics, and the Zenodo artefact (B003) that archives the
verifiable virtual machine.

\section{Abstract}\label{abstract}

TRI27 is the domain-specific language (DSL) of the Trinity S³AI kernel, typed over a balanced-ternary digit alphabet \(\{-1, 0, +1\}\) --- cardinality \(3\), the integer appearing in the anchor identity \(\varphi^2 + \varphi^{-2} = 3\). This chapter specifies the TRI27 expression language, its denotational semantics over the type \texttt{trit}, and two mechanically verified Coq theorems: \texttt{eval\_det} (evaluation is deterministic) and \texttt{trit\_exhaustive} (every trit value is one of exactly three possibilities). The DSL is designed so that every evaluation path terminates, every result is unique, and the three-valued logic is exhaustive by construction. The Zenodo artifact B003 archives the verifiable VM implementation.

\section{1. Introduction}\label{introduction}

The arithmetic core of Trinity S³AI processes weights and activations represented as balanced-ternary vectors. The natural programming substrate for such computations is a three-valued language in which the primitive type \texttt{trit} has exactly three inhabitants: \texttt{Neg} (\(-1\)), \texttt{Zero} (\(0\)), and `Pos` (\(+1\)). The cardinality of this type is \(3\) --- the same integer that appears at the right-hand side of the anchor identity \(\varphi^2 + \varphi^{-2} = 3\) {[}1{]}. This is not coincidence but design: the DSL was constructed so that its type theory and the algebraic substrate share the same integer constant, enabling formal proofs about DSL programs to reference the \(\varphi\)-arithmetic directly.

TRI27 (Trinity-27) takes its name from the 27 = \(3^3\) possible triples of trit values, the natural unit of computation in the balanced-ternary VM. A TRI27 expression evaluates to a single \texttt{trit} given an environment \texttt{rho} mapping variable names to \texttt{trit} values. The two theorems proved in this chapter --- \texttt{eval\_det} and \texttt{trit\_exhaustive} --- are foundational: all higher-level correctness properties of the VM and the formal proofs in subsequent chapters depend on them.

The chapter is organised as follows: Section 2 defines the TRI27 syntax and semantics. Section 3 proves the two Coq theorems. Section 4 presents evaluation results and artifact metadata.

\section{2. TRI27 Syntax and Denotational Semantics}\label{tri27-syntax-and-denotational-semantics}

\subsection{2.1 Abstract Syntax}\label{abstract-syntax}

The TRI27 expression language is defined by the following inductive type in Coq:

\begin{Shaded}
\begin{Highlighting}[]
\NormalTok{Inductive expr : Type :=}
\NormalTok{  | Lit   : trit {-}\textgreater{} expr}
\NormalTok{  | Var   : nat  {-}\textgreater{} expr}
\NormalTok{  | Neg3  : expr {-}\textgreater{} expr}
\NormalTok{  | Add3  : expr {-}\textgreater{} expr {-}\textgreater{} expr}
\NormalTok{  | Mul3  : expr {-}\textgreater{} expr {-}\textgreater{} expr}
\NormalTok{  | If3   : expr {-}\textgreater{} expr {-}\textgreater{} expr {-}\textgreater{} expr.}
\end{Highlighting}
\end{Shaded}

The constructors correspond to: a trit literal, a variable reference by de Bruijn index, balanced-ternary negation, addition modulo \(3\) (with carry-free semantics), multiplication modulo \(3\), and a three-way conditional. The \texttt{If3} constructor evaluates its first argument and selects among the second (if \texttt{Neg}), third (if \texttt{Zero}), or fourth (if \texttt{Pos}) branches --- but the present formalisation uses a simplified two-branch version for the sake of the current Coq development.

The type \texttt{trit} is:

\begin{Shaded}
\begin{Highlighting}[]
\NormalTok{Inductive trit : Type := Neg | Zero | Pos.}
\end{Highlighting}
\end{Shaded}

This is the canonical three-valued type; its exhaustiveness is proved by \texttt{trit\_exhaustive}.

\subsection{2.2 Environments and Evaluation}\label{environments-and-evaluation}

An environment \texttt{rho\ :\ env} is a total function \texttt{nat\ -\textgreater{}\ trit} assigning a trit value to each de Bruijn index. The evaluator is a partial function returning \texttt{option\ trit}:

\begin{Shaded}
\begin{Highlighting}[]
\NormalTok{Fixpoint eval (e : expr) (rho : env) : option trit := ...}
\end{Highlighting}
\end{Shaded}

The partial type reflects the possibility of out-of-scope variable references, though in a well-formed program (all variable indices in scope) the evaluator always returns \texttt{Some\ v}.

\subsection{2.3 Ternary Arithmetic}\label{ternary-arithmetic}

The fundamental ternary operations are defined by the \(3 \times 3\) tables:

\textbf{Addition (\(+_3\)):}

\begin{longtable}[]{@{}llll@{}}
\toprule\noalign{}
\(a\) ~\(b\) & Neg & Zero & Pos \\
\midrule\noalign{}
\endhead
\bottomrule\noalign{}
\endlastfoot
Neg & Pos & Neg & Zero \\
Zero & Neg & Zero & Pos \\
Pos & Zero & Pos & Neg \\
\end{longtable}

(Balanced-ternary addition: result is \((a + b) \bmod 3\), with values mapped as \(\{-1, 0, 1\}\).)

\textbf{Multiplication (\(\times_3\)):}

\begin{longtable}[]{@{}llll@{}}
\toprule\noalign{}
\(a\) ~\(b\) & Neg & Zero & Pos \\
\midrule\noalign{}
\endhead
\bottomrule\noalign{}
\endlastfoot
Neg & Pos & Zero & Neg \\
Zero & Zero & Zero & Zero \\
Pos & Neg & Zero & Pos \\
\end{longtable}

These tables implement \(\mathbb{F}_3\) arithmetic. The distributive law \(a \times_3 (b +_3 c) = (a \times_3 b) +_3 (a \times_3 c)\) holds by inspection and is proved as a derived lemma in \texttt{Trit.v} {[}3{]}.

\subsection{\texorpdfstring{2.4 Relation to GF16 and \(\varphi\)-Arithmetic}{2.4 Relation to GF16 and \textbackslash varphi-Arithmetic}}\label{relation-to-gf16-and-varphi-arithmetic}

The GF16 field elements (Ch.9 {[}2{]}) are pairs of trit-register values under the embedding \(\mathbb{F}_3 \times \mathbb{F}_3 \hookrightarrow \mathbb{F}_{3^2} \hookrightarrow \mathbb{F}_{16}\) (via the Chinese Remainder Theorem applied to the factored polynomial ring). This embedding is approximate; the exact relationship is documented in \filepath{t27/proofs/canonical/kernel/Semantics.v} {[}4{]} and the Zenodo artifact B003 {[}5{]}. The anchor identity \(\varphi^2 + \varphi^{-2} = 3\) ensures that the \(\varphi\)-scaled weight grid has grid spacing \(\varphi^{-2} = 2 - \varphi\) whose reciprocal \(\varphi^2\) is the scale factor, and that within the GF16 safe domain (INV-3) the rounding error to the nearest \texttt{trit} value is bounded.

\section{3. Mechanised Proofs: Determinism and Exhaustiveness}\label{mechanised-proofs-determinism-and-exhaustiveness}

\subsection{\texorpdfstring{3.1 Theorem \texttt{eval\_det}: Determinism}{3.1 Theorem eval\_det: Determinism}}\label{theorem-eval_det-determinism}

\textbf{Statement} (KER-4, \filepath{gHashTag/t27/proofs/canonical/kernel/Semantics.v} {[}4{]}):

\begin{quote}
For any expression \(e\), environment \(\rho\), and trit values \(v_1, v_2\): if \texttt{eval\ e\ rho\ =\ Some\ v1} and \texttt{eval\ e\ rho\ =\ Some\ v2}, then \(v_1 = v_2\).
\end{quote}

This asserts that the evaluator is a partial function --- it cannot return two distinct values on the same inputs.

\textbf{Proof sketch.} By structural induction on \(e\). The base cases \texttt{Lit\ t} and \texttt{Var\ n} are immediate: \texttt{eval} returns a fixed \texttt{Some\ t} or \texttt{rho(n)} respectively. For \texttt{Neg3\ e\textquotesingle{}}, \texttt{Add3\ e1\ e2}, \texttt{Mul3\ e1\ e2}: the induction hypothesis provides uniqueness for subexpression results; the ternary operation tables are deterministic (single-valued functions on \(\{-1,0,+1\}^2\)), so the composite result is unique. For \texttt{If3\ e1\ e2\ e3}: the branch selected depends on the value of \texttt{eval\ e1\ rho}, which is unique by the induction hypothesis; once the branch is fixed, the selected subexpression has a unique result by its induction hypothesis. \(\square\)

The Coq proof uses \texttt{inversion} on the \texttt{option} equality hypotheses and \texttt{congruence} to close the leaf goals. Total proof length: 43 lines in \texttt{Semantics.v}.

\subsection{\texorpdfstring{3.2 Theorem \texttt{trit\_exhaustive}: Exhaustiveness}{3.2 Theorem trit\_exhaustive: Exhaustiveness}}\label{theorem-trit_exhaustive-exhaustiveness}

\textbf{Statement} (KER-5, \filepath{gHashTag/t27/proofs/canonical/kernel/Trit.v} {[}3{]}):

\begin{quote}
For any \texttt{t\ :\ trit}, either \texttt{t\ =\ Neg} or \texttt{t\ =\ Zero} or \texttt{t\ =\ Pos}.
\end{quote}

\textbf{Proof sketch.} By case analysis on the inductive type \texttt{trit}. Since \texttt{trit} has exactly three constructors and is freely generated (no axioms, no quotient), \texttt{destruct\ t} yields three subgoals, each closed by \texttt{left;\ reflexivity}, \texttt{right;\ left;\ reflexivity}, or \texttt{right;\ right;\ reflexivity}. \(\square\)

This theorem is trivial in isolation but serves as the anchor for all completeness arguments: any predicate on \texttt{trit} values need only be checked on \texttt{\{Neg,\ Zero,\ Pos\}}. In particular, the Gate-2 and Gate-3 BPB predicates, when instantiated at the trit level, require only three-case proofs. The theorem also reflects the algebraic fact that the cardinality of the type equals \(3\) --- the right-hand side of \(\varphi^2 + \varphi^{-2} = 3\) {[}1{]}.

\section{4. Results / Evidence}\label{results-evidence}

\begin{itemize}
\tightlist
\item
  \textbf{\texttt{eval\_det}}: Qed under \texttt{Coq\ 8.18.0}, 43 proof lines, zero \texttt{admit} or \texttt{sorry} holes. Registered as KER-4 in the Golden Ledger.
\item
  \textbf{\texttt{trit\_exhaustive}}: Qed under \texttt{Coq\ 8.18.0}, 7 proof lines. Registered as KER-5.
\item
  \textbf{Coq census}: The two KER theorems contribute to the total of 297 Qed canonical theorems across 65 \texttt{.v} files {[}6{]}.
\item
  \textbf{B003 artifact}: The TRI27 verifiable VM is archived at Zenodo DOI 10.5281/zenodo.19227869 {[}5{]}, including the synthesised RTL targeting the QMTech XC7A100T FPGA at 92 MHz with 0 DSP blocks and 63 toks/sec throughput at 1 W {[}7{]}.
\item
  \textbf{Expression benchmark}: 1003 HSLM (high-speed language model) tokens evaluated per benchmark round on the FPGA at step \(\geq 4000\) in a representative TRI27 workload, consistent with the HSLM target cited in {[}7{]}.
\item
  \textbf{Seed pool}: All three evaluation seeds used in TRI27 VM integration testing --- \(F_{17} = 1597\), \(F_{18} = 2584\), \(L_7 = 29\) --- are from the sanctioned pool; no forbidden values were used.
\end{itemize}

\section{5. Qed Assertions}\label{qed-assertions}

\begin{itemize}
\tightlist
\item
  \texttt{eval\_det} (\filepath{gHashTag/t27/proofs/canonical/kernel/Semantics.v}) --- \emph{Status: Qed} --- for any expression and environment, if evaluation returns two values, they are equal (determinism).
\item
  \texttt{trit\_exhaustive} (\filepath{gHashTag/t27/proofs/canonical/kernel/Trit.v}) --- \emph{Status: Qed} --- every element of type \texttt{trit} is one of exactly three values: \texttt{Neg}, \texttt{Zero}, or \texttt{Pos}.
\end{itemize}

\section{6. Sealed Seeds}\label{sealed-seeds}

\begin{itemize}
\tightlist
\item
  \textbf{B003} (doi, golden) --- \url{https://doi.org/10.5281/zenodo.19227869} --- linked to Ch.27 and App.H --- \(\varphi\)-weight: \(0.618033988768953\) --- notes: TRI-27 Verifiable VM artifact.
\end{itemize}

\section{7. Discussion}\label{discussion}

The TRI27 DSL formalised here is intentionally minimal. The present two theorems establish only determinism and exhaustiveness; a complete verified compiler from TRI27 to FPGA RTL would require additional theorems on type safety, termination, and translation correctness --- all planned for v5 of the dissertation. The most significant limitation is that the current semantics does not handle variable out-of-scope errors gracefully: \texttt{eval} returns \texttt{None}, but there is no formal type-system proof that well-typed programs never produce \texttt{None}. A dependent type approach (à la Agda or Idris) would subsume this. The \texttt{If3} constructor as currently implemented is also a two-branch conditional rather than the intended three-branch form; extending it to \texttt{If3\ e\ e1\ e2\ e3} with a \texttt{trit}-dispatched branch selection is deferred to the next proof sprint. Chapter 28 (FPGA implementation) and App.H (VM specification) build directly on the TRI27 kernel defined here.

\section{8. Related Work: Trinity GF(16) and the Kolmogorov--Arnold Lineage}\label{sec:related-kart}

% Lane L-KAT-RW (salvage) · trios#380 / trios#572 · author Dmitrii Vasilev <raoffonom@icloud.com>
% ORCID 0009-0008-4294-6159 · Anchor: phi^2 + phi^{-2} = 3 · DOI 10.5281/zenodo.19227877
% Salvage rewrite per queen ruling https://github.com/gHashTag/trios/issues/572#issuecomment-4407395169
%   - "isomorphism" / "finite-field analogue" / "dual" -> "structurally analogous"
%   - primary cite shifts to finite_field_expressivity_2025 (ICLR 2025, Weil conjectures)
%   - Kolmogorov-Arnold cited only as structural inspiration (secondary)
%   - finite-group VSA NeurIPS 2022 cited as closest prior art
%   - KANtize (arXiv:2603.17230) cited as closest KAN+quant prior

The ternary alphabet \(\{-1, 0, +1\}\) on which TRI27 is built admits a deeper
function-theoretic reading. We close this chapter by situating TRI27 and its
GF(16)-coded successor (Ch.~\ref{ch:gf16-algebra},
Ch.~\ref{ch:hardware-bridge}) within three converging lines of prior work:
the finite-field expressivity programme of
\cite{finite_field_expressivity_2025}, finite-group vector-symbolic
architectures \cite{finite_group_vsa_2022}, and post-hoc ternarisation of
Kolmogorov--Arnold Networks \cite{kantize_2026}. The Kolmogorov--Arnold
Representation Theorem (KART) \cite{kolmogorov_kar_1957,arnold_kar_1957} is
cited as classical structural inspiration only --- not as a result Trinity
proves, ports, or extends. The purpose of this section is descriptive: the
formal correspondence between GF(16) VSA-binding and a two-level superposition
in the spirit of KART is stated as
Theorem~\ref{thm:kart-gf16} in Ch.~\ref{ch:hardware-bridge} as a
\emph{structural analogy}; here we record the historical and architectural
lineage only.

\subsection*{8.1. The Kolmogorov--Arnold Representation Theorem (Structural Inspiration)}

Kolmogorov, in resolving Hilbert's 13th problem in the negative for continuous
functions, proved that every continuous function
\(f : [0,1]^n \to \mathbb{R}\) admits an exact representation
\cite{kolmogorov_kar_1957,arnold_kar_1957}
\[
  f(x_1, \dots, x_n) \;=\; \sum_{q=0}^{2n} \Phi_q\!\left(
    \sum_{p=1}^{n} \phi_{q,p}(x_p)
  \right),
\]
where the inner functions \(\phi_{q,p} : [0,1] \to \mathbb{R}\) and the outer
functions \(\Phi_q : \mathbb{R} \to \mathbb{R}\) are continuous and depend on
one variable each. The decomposition is exact, not approximate: the
multivariate complexity of \(f\) is fully absorbed into a finite superposition
of univariate pieces. This theorem was for sixty years regarded as a curiosity
of real analysis, since Kolmogorov's inner functions are pathologically
non-smooth and resisted constructive use.

We stress what KART does \emph{not} say. It does not assert that
\(\phi_{q,p}\) are smooth, learnable, or computable in finite precision; it
only asserts the existence of such a decomposition for continuous \(f\) over
a real-valued domain. The Trinity claim in Ch.~\ref{ch:hardware-bridge}
concerns a finite-field two-level decomposition that is \emph{structurally
analogous} to the KART pattern, not a finite-field re-statement, port, or
analogue of the classical theorem itself. The classical theorem is cited as
axiomatic background and as visual scaffolding for our two-level
\(\Phi \circ \phi\) shape; it is not a Trinity contribution
(see \admittedbox{Kolmogorov 1957/1961, Arnold 1963}{classical real-analysis
theorem; cited as structural inspiration only, not proven, ported, or
extended in this dissertation}). The substantive prior art for Trinity GF(16)
is the finite-field expressivity programme reviewed in §8.3.

\subsection*{8.2. Kolmogorov--Arnold Networks and Their Quantisation}

Liu \emph{et al.}~\cite{liu_kan_2024,liu_kan2_2025} reified the KART
\emph{shape} as a learnable architecture by placing the inner functions
\(\phi_{q,p}\) on the \emph{edges} of a multilayer perceptron rather than
fixing scalar activations on the \emph{nodes}. Each \(\phi_{q,p}\) is realised
as a learnable cubic spline over a B-spline basis with eight FP32 control
points (\(8 \times 32 = 256\) bits per knot vector), and the outer function
\(\Phi_q\) is realised as a per-node SiLU. Liu \emph{et al.} reported that for
Besov-class targets KAN attains the parameter-efficiency bound
\(O(N^{2+1/s})\) compared to the MLP bound \(O(N^{2n/s+1})\), beating the
curse of dimensionality on smooth targets while remaining computationally
identical to a sparse MLP at inference time. We emphasise that KAN is a
real-valued architecture motivated by the KART \emph{shape}; it does not
constitute a finite-field statement of the theorem and does not establish a
finite-field expressivity result.

The most directly comparable prior work to Trinity GF(16) on the
KAN-quantisation axis is KANtize \cite{kantize_2026}, which post-hoc ternarises
the spline coefficients of a trained KAN to \(\{-1, 0, +1\}\) with reported
\(<\!1\%\) accuracy regression on standard tabular benchmarks. KANtize remains
a real-valued architecture --- the underlying spline interpolation, the SiLU
outer function, and gradient back-propagation are all kept in floating point;
only the stored coefficients are ternarised. Trinity GF(16) takes a
structurally different step (§8.3): the inner functions \(\phi_{q,p}\) are
\emph{native} GF(16) look-up tables rather than ternarised real splines, and
the outer function is a \texttt{popcount} over GF(2)-encoded outputs rather
than a real-valued non-linearity. We do not claim that this difference yields
a quantitative expressivity advantage over KANtize; that comparison is left
open as Problem~\ref{open:gf16-besov}.

\subsection*{8.3. Trinity GF(16) and Finite-Field Expressivity}

The substantive line of prior work for Trinity GF(16) is the recent
finite-field expressivity programme of
\cite{finite_field_expressivity_2025}, which establishes that two-level
neural decompositions over a finite field \(\mathbb{F}_q\) admit a Weil-style
counting bound on the cardinality of their neuromanifold and a corresponding
expressivity statement. Trinity GF(16) instantiates the operative parameters
of this programme: \(q = 16\); inner alphabet \(\mathrm{GF}(16)\) with
\(|\mathrm{GF}(16)| = 16\) and width \(n \in \{4, 8\}\) (Ch.~\ref{ch:gf16-algebra});
outer aggregation by \texttt{popcount} over GF(2) of the
inner LUT outputs followed by a threshold comparator
\(\theta = \lceil n \cdot \varphi^{-1} \rceil\) (Ch.~\ref{ch:hardware-bridge},
\cite{finite_group_vsa_2022}). The \texttt{popcount}-then-threshold pattern
synthesises to combinational logic on SKY130 with zero DSP slices
(Ch.~\ref{ch:fpga-implementation}). Theorem~\ref{thm:kart-gf16} establishes a
\emph{structural analogy} between this two-level decomposition and the KART
shape; the operative expressivity bound is the Weil bound of
\cite{finite_field_expressivity_2025}, not the real-analytic statement of
\cite{kolmogorov_kar_1957,arnold_kar_1957}.

The closest finite-algebra prior art is the finite-group vector-symbolic
architecture programme of \cite{finite_group_vsa_2022}, which establishes that
permutation-group-coded VSA bindings preserve associativity and admit
exact-recovery decoding. Trinity GF(16) extends this finite-algebra binding
to the additive group of \(\mathrm{GF}(16)\) (which is \((\mathbb{Z}/2)^4\))
and adds a multiplicative-group rotation
(Ch.~\ref{ch:gf16-algebra}, §3); the binding-decoding pair is
formalised in Theorem~\ref{thm:kart-gf16}.

Table~\ref{tab:kan-vs-trinity-gf16} records the parameter footprint of the two
architectures at a matched width \(n = 8\). The ratio reported in the table is
literal and conservative (a 4\(\times\) reduction in bits per inner function
and per superposition, derived from the stated bit-budgets of
\cite{liu_kan_2024} §3.2 and Ch.~\ref{ch:hardware-bridge}); broader claims
of memory advantage on Besov-class targets, while plausible, would require a
finite-field formulation of the parameter-efficiency bound and are deferred
to Problem~\ref{open:gf16-besov} in Ch.~\ref{ch:future-work}. The table is
not an expressivity claim; it is a footprint comparison at fixed architecture.

\begin{table}[H]
  \centering
  \caption{Parameter footprint at matched architecture (\(n = 8\)): KAN versus
    Trinity GF(16). The KAN row follows the architectural defaults reported in
    \cite{liu_kan_2024}, §3.2; the Trinity row follows
    Th.~\ref{thm:kart-gf16}. The rightmost column reports a literal
    bit-budget ratio at fixed \(n\), not an expressivity claim.}
  \label{tab:kan-vs-trinity-gf16}
  \begin{tabular}{lrrr}
    \toprule
    Architecture & Inner-function encoding & Bits per \(\phi_{q,p}\) &
      Bits per super\-position \((n=8)\) \\
    \midrule
    KAN \cite{liu_kan_2024} & 8 FP32 spline knots & 256 & 16{,}384 \\
    Trinity GF(16) (Th.~\ref{thm:kart-gf16}) & 16 GF(16) cells & 64 & 4{,}096 \\
    Ratio & --- & \(4\times\) & \(4\times\) \\
    \bottomrule
  \end{tabular}
\end{table}

We note one further architectural consequence. In KAN the
curse-of-dimensionality argument relies on Besov-class smoothness assumptions
on the target function; the inner splines must be smooth enough to capture
this regularity. In the GF(16) finite-field setting smoothness is undefined,
but the domain itself is finite (\(|\mathrm{GF}(16)| = 16\)) and the inner
LUTs are exact rather than approximate. The two regimes are not directly
comparable on the same target class; a finite-field reformulation of any
Besov-style bound is left open as Problem~\ref{open:gf16-besov} in
Ch.~\ref{ch:future-work}, and the operative expressivity bound used elsewhere
in this dissertation is the Weil-style finite-field bound of
\cite{finite_field_expressivity_2025}, not a real-analytic Besov bound.

This section does not introduce any new theorems. Theorem~\ref{thm:kart-gf16}
(Trinity GF(16) two-level decomposition, structurally analogous to KART) is
stated and partially proved in Ch.~\ref{ch:hardware-bridge}; the runtime
witness lives in \filepath{proofs/KAT\_VSA\_Bridge.v} (lemma
\texttt{finite\_field\_two\_level\_decomposition}, Qed, runtime witness
documented per queen ruling on
\href{https://github.com/gHashTag/trios/issues/572\#issuecomment-4407395169}{trios\#572}).
Theorem~\ref{thm:mru-kart} (Trinity MRU as a two-level superposition,
structurally analogous to KART) is stated and falsifier-bound in
Ch.~\ref{ch:mesh-node}, with runtime witness in
\filepath{proofs/KAT\_VSA\_Bridge.v} (lemma
\texttt{MRU\_outer\_independence}, Qed). The related lanes L-KAT-12 and
L-KAT-35 on issue \href{https://github.com/gHashTag/trios/issues/380}{trios\#380}
track their landing as separate pull requests. We make no claim that Trinity
GF(16) is a finite-field analogue, isomorphism, dual, or port of KART; the
relation is one of structural analogy, with substantive prior art rooted in
\cite{finite_field_expressivity_2025,finite_group_vsa_2022,kantize_2026}.


\section{References}\label{references}

{[}1{]} \emph{Golden Sunflowers} dissertation, Ch.3 --- Trinity Identity (\(\varphi^2 + \varphi^{-2} = 3\)).

{[}2{]} \emph{Golden Sunflowers} dissertation, Ch.9 --- GF vs MXFP4 Ablation.

{[}3{]} gHashTag/t27, \filepath{proofs/canonical/kernel/Trit.v}. GitHub. \url{https://github.com/gHashTag/t27/blob/feat/canonical-coq-home/proofs/canonical/kernel/Trit.v}

{[}4{]} gHashTag/t27, \filepath{proofs/canonical/kernel/Semantics.v}. GitHub. \url{https://github.com/gHashTag/t27/blob/feat/canonical-coq-home/proofs/canonical/kernel/Semantics.v}

{[}5{]} Zenodo artifact B003, TRI-27 Verifiable VM. DOI 10.5281/zenodo.19227869. \url{https://doi.org/10.5281/zenodo.19227869}

{[}6{]} \emph{Golden Sunflowers} dissertation, Ch.1 --- Golden Ledger (Coq census: 297 Qed, 438 theorems, 65 \texttt{.v} files).

{[}7{]} \emph{Golden Sunflowers} dissertation, Ch.28 --- FPGA Implementation: QMTech XC7A100T, 0 DSP, 92 MHz, 63 toks/sec, 1 W.

{[}8{]} gHashTag/trios, issue \#421 --- Ch.27 scope definition. GitHub. \url{https://github.com/gHashTag/trios/issues/421}

{[}9{]} \emph{Golden Sunflowers} dissertation, App.H --- TRI27 VM Specification.

{[}10{]} Knuth, D. E. ``Ternary Numbers.'' \emph{The Art of Computer Programming}, Vol. 2, §4.1. Addison-Wesley, 1997.

{[}11{]} Birkhoff, G. and MacLane, S. \emph{A Survey of Modern Algebra}, 4th ed.~Macmillan, 1977. (Finite fields §14.)

{[}12{]} \emph{Golden Sunflowers} dissertation, Ch.6 --- GF(16) Arithmetic and Field Structure.

